\section{Technology Stack}

The implementation of SIGS-RS leverages modern asynchronous runtimes, time-series storage engines, distributed key-value caches, and cloud-native container orchestration. Every technology selected for the platform was audited against the authoritative SRS non-functional criteria (\texttt{REQ-PF-1}, \texttt{REQ-PF-2}, \texttt{REQ-PF-3}, and \texttt{REQ-SE-1} through \texttt{REQ-SE-4}). Table~\ref{tab:tech_stack} enumerates the confirmed technology components across each layer of the gateway architecture.

\begin{table}[H]
\centering
\caption{SIGS-RS Verified Technology Stack and Runtime Dependencies}
\label{tab:tech_stack}
\small
\begin{tabularx}{\linewidth}{@{}p{2.8cm} p{2.8cm} p{1.8cm} X@{}}
\toprule
\textbf{Subsystem Layer} & \textbf{Technology / Library} & \textbf{Version} & \textbf{Architectural Responsibility} \\
\midrule
\textbf{Language Runtime} & Python & 3.12.3 & Core runtime providing native asynchronous I/O (\texttt{asyncio}), structured task groups, and memory-safe garbage collection. \\
\textbf{Web / API Framework} & FastAPI & 0.115.0 & High-performance asynchronous REST API framework providing automatic OpenAPI schema generation, dependency injection, and sub-millisecond route dispatch. \\
\textbf{ASGI Server} & Uvicorn & 0.30.0 & Lightning-fast ASGI web server built on \texttt{uvloop} and \texttt{httptools}, executing the gateway event loop. \\
\textbf{Schema Validation} & Pydantic & 2.8.2 & Rust-accelerated schema validation and serialization engine, enforcing strict telemetry payload types and impossible-value validation. \\
\textbf{MQTT Ingress} & aiomqtt & 2.1.0 & Asynchronous MQTT client supporting TLS 1.2+ mutual authentication, topic wildcard subscriptions (\texttt{telemetry/\#}), and non-blocking packet ingestion. \\
\textbf{Rate-Limit Engine} & Redis / aioredis & 7.2.4 & In-memory distributed datastore providing atomic sliding-window transaction scripts, sorted set expiration, and cluster synchronization. \\
\textbf{Time-Series Persistence} & TimescaleDB / PostgreSQL & 16.2 / 2.14 & Relational time-series database managing partitioned hypertables for sanitized readings and relational integrity for device profiles. \\
\textbf{Database Driver} & asyncpg & 0.29.0 & Native asynchronous PostgreSQL protocol client offering binary protocol encoding and high-throughput connection pooling. \\
\textbf{Authentication} & PyJWT \& cryptography & 2.8.0 & Cryptographic signature verification for OAuth 2.0 / JWT tokens with role-based claim extraction (\texttt{security\_engineer}, \texttt{compliance\_auditor}). \\
\textbf{Metrics Exposition} & prometheus\_client & 0.20.0 & In-process metric collector exposing standard Prometheus counters, gauges, and latency histograms at \texttt{/metrics}. \\
\textbf{Testing Framework} & pytest \& pytest-asyncio & 8.3.2 & Test execution runner driving 580 automated test cases across unit, integration, and microbenchmark suites. \\
\textbf{Containerization} & Docker & 26.1 & Multi-stage container builds generating minimal, distroless production images for the gateway process. \\
\textbf{Orchestration} & Kubernetes & 1.30 & Cloud-native workload management featuring Deployment manifests, Horizontal Pod Autoscalers, ConfigMaps, and Secret management. \\
\bottomrule
\end{tabularx}
\end{table}

\subsection{Language and Asynchronous Runtime Rationale}
Python 3.12 was selected as the foundational runtime for SIGS-RS due to its mature asynchronous I/O framework (\texttt{asyncio}), which enables concurrent handling of thousands of open network sockets without the memory overhead of OS-level OS threading. In classical multi-threaded architectures, sustaining 10,000 concurrent sensor connections incurs significant context-switching penalties and memory allocation ($\sim 8$~MB stack per thread). In contrast, Python's cooperative multitasking model schedules lightweight coroutines onto a single-threaded event loop, consuming less than 1~KB of state per connection.

Furthermore, Python 3.12 introduces substantial interpreter optimizations, including adaptive specialization (PEP 659), faster exception handling, and reduced object allocation overhead, which are instrumental in satisfying the strict 15.0-millisecond latency ceiling mandated by \texttt{REQ-PF-2}.

\subsection{Database and Cache Selection}
The gateway storage architecture strictly isolates high-throughput sliding-window accounting from permanent time-series persistence:
\begin{enumerate}
    \item \textbf{Redis (Sliding Window):} Evaluated against in-memory hash tables, Redis Sorted Sets (\texttt{ZSET}) provide atomic scoring by Unix microsecond timestamp. Using Lua script pipelining, the inspector records arrivals and trims entries older than $T - 10$ seconds in a single round-trip, executing in under 0.2~ms.
    \item \textbf{TimescaleDB (Hypertable Storage):} Standard relational databases suffer from index tree degradation as table size grows past millions of rows. TimescaleDB solves this through automatic two-dimensional partitioning into immutable time chunks (\textit{hypertables}), preserving sustained write throughput of tens of thousands of inserts per second while supporting full SQL querying for analytical dashboards.
\end{enumerate}
